% !TEX root = ../Main.tex
\chapter{放回与不放回}

抽取问题的两大模式：\textbf{有放回}（每次取完放回原处，下次仍可能取到该物）与\textbf{无放回}（取完不放回，下次样本库少一个）。两者的\textbf{样本空间大小不同}，但在对称性下\textbf{概率值往往相同}——这是一个微妙的直觉陷阱。

\section{不放回抽取}

\state{不放回的样本空间}{
    有 $N$ 个物品，不放回地抽取 $n$ 次，每次 $1$ 个（$n \le N$）：
    \begin{itemize}
        \item \textbf{有顺序}：排列 $A_N^n = \dfrac{N!}{(N-n)!}$；
        \item \textbf{无顺序}：组合 $C_N^n = \dfrac{N!}{n!(N-n)!}$。
    \end{itemize}

    \textbf{如何选样本空间}：
    \begin{itemize}
        \item 题目问"恰有、至少、同色"等\textbf{只关心哪些物品被取到}——用组合；
        \item 题目问"第一次是……第二次是……"等\textbf{关心顺序}——用排列。
    \end{itemize}
}

\ex[不放回取两数]{
    在 $0, 1, 2, 3, 4, 5, 6, 7, 8$ 这 $9$ 个数中\textbf{不放回}地取两个。记事件：
    \begin{itemize}
        \item $A$ = "取出的两个数中\textbf{有且仅有一个}为奇数"；
        \item $B$ = "取出的两个数中\textbf{有且仅有一个}为偶数"。
    \end{itemize}
    问是否有 $P(A) = P(B)$？
}

\sol{
    \textbf{关键观察}：两次抽取共取 $2$ 个数。"恰有一个奇数" $\iff$ "恰有一个偶数"（$2 - 1 = 1$）。故\textbf{事件 $A$ 与 $B$ 本身相同}，$P(A) = P(B)$ 必成立。

    \textbf{验证数值}：$9$ 个数中奇数 $4$ 个（$\{1,3,5,7\}$），偶数 $5$ 个（$\{0,2,4,6,8\}$）。
    \[
    P(A) = \frac{C_4^1 \cdot C_5^1}{C_9^2} = \frac{4 \cdot 5}{36} = \frac{5}{9},\qquad P(B) = \frac{C_5^1 \cdot C_4^1}{C_9^2} = \frac{5 \cdot 4}{36} = \frac{5}{9}.
    \]
    确实相等。$\checkmark$
}

\section{有放回抽取}

\state{有放回的样本空间}{
    有 $N$ 个物品，有放回地抽取 $n$ 次，每次 $1$ 个——\textbf{每次取完放回原处}。
    \begin{itemize}
        \item \textbf{样本空间大小}：$N^n$（每次独立有 $N$ 种选择）；
        \item \textbf{顺序是"内置"的}——有放回时通常天然带顺序（第一次、第二次……）。
    \end{itemize}
}

\ex[例~1 的有放回变形]{
    把上例改为"\textbf{有放回}地取两次"，其它条件不变。问是否有 $P(A) = P(B)$？
}

\sol{
    样本空间为有序对，$|\Omega| = 9^2 = 81$。

    $A$ = "恰有一个奇数" 的样本点数：
    \begin{itemize}
        \item 第一次奇（$4$ 种）、第二次偶（$5$ 种）：$4 \cdot 5 = 20$；
        \item 第一次偶（$5$ 种）、第二次奇（$4$ 种）：$5 \cdot 4 = 20$。
    \end{itemize}
    合计 $40$。$P(A) = \dfrac{40}{81}$。同理 $P(B) = \dfrac{40}{81}$。

    \textbf{$P(A) = P(B)$ 仍成立。}$\checkmark$ 原因和例 1 一样——\textbf{"恰一奇" 与 "恰一偶" 在两次抽取中逻辑等价}，与是否放回无关。
}

\rmkb{
    \textbf{重要观察}：若两事件\textbf{本身逻辑等价}（$A = B$ 作为集合），概率必相等——与放回方式、样本空间构造都无关，只与事件的\textbf{逻辑定义}有关。本题的"坑"在于题目刻意把同一事件写成两条看似不同的语句。
}

\section{对比：放回 vs 不放回的典型题}

\ex[同色球 / 至少一个]{
    袋中有 $7$ 个红球、$3$ 个绿球。\textbf{无放回}地依次抽取 $2$ 次，每次 $1$ 个。求：
    \begin{enumerate}
        \item 取得 $2$ 个红球的概率；
        \item 取得 $2$ 个绿球的概率；
        \item 取得 $2$ 个同色球的概率；
        \item 至少取得 $1$ 个红球的概率。
    \end{enumerate}
}

\sol{
    共 $10$ 个球，不放回取 $2$ 次（有序）的样本空间大小 $10 \cdot 9 = 90$。

    \textbf{(1)} 两红球的顺序对数 $7 \cdot 6 = 42$。$P = \dfrac{42}{90} = \dfrac{7}{15}$。

    \textbf{(2)} 两绿球 $3 \cdot 2 = 6$。$P = \dfrac{6}{90} = \dfrac{1}{15}$。

    \textbf{(3)} 同色 = 两红 或 两绿（互斥）：
    \[
    P = \frac{7}{15} + \frac{1}{15} = \frac{8}{15}.
    \]

    \textbf{(4)} "至少一个红" $\iff$ "并非全为绿"，对立事件为"两个都是绿球"：
    \[
    P = 1 - \frac{1}{15} = \frac{14}{15}.
    \]
}

\rmkb{
    \textbf{古典概型解题三步曲}：
    \begin{enumerate}
        \item \textbf{刻画样本空间}：列出等可能的样本点——\textbf{有序}或\textbf{无序}要与事件描述一致；
        \item \textbf{刻画事件}：数出事件对应的样本点个数；
        \item \textbf{套公式}：$P(A) = |A| / |\Omega|$。
    \end{enumerate}
    \textbf{最常见的坑}——"至少"类问题：若事件可以"又……又……"组合，\textbf{对立事件}往往更简单。
}
